\lampiransection{Format Survei Pelanggan Pendukung}
\label{app:survei-pelanggan-pendukung}

Lampiran ini digunakan sebagai format survei pelanggan pendukung. Survei tidak diposisikan sebagai analisis statistik utama, tetapi sebagai data tambahan untuk membaca pengalaman pelanggan, alasan membeli, peluang pembelian ulang, dan temuan observasi GrabFood yang perlu dikonfirmasi dari sisi pelanggan.

\begin{table}[!htbp]
    \centering
    \caption{Rancangan Pertanyaan Survei Pelanggan Setelah Observasi GrabFood}
    \label{tab:rancangan-survei-pelanggan}
    \scriptsize
    \setlength{\tabcolsep}{3pt}
    \renewcommand{\arraystretch}{0.92}
    \begin{tabularx}{\textwidth}{|Z{0.8cm}|Y|Z{2.6cm}|Y|}
        \toprule
        \textbf{No.} & \textbf{Pertanyaan} & \textbf{Bentuk Jawaban} & \textbf{Tujuan Data} \\ \hline
        \midrule
        1 & Apakah pernah membeli Nasi Gerilya melalui GrabFood? & Ya / Tidak & Penyaring responden agar jawaban berasal dari pengalaman aktual. \\ \hline
        2 & Kapan terakhir membeli Nasi Gerilya melalui GrabFood? & Pilihan periode & Membedakan pengalaman terbaru dan pengalaman lama. \\ \hline
        3 & Berapa kali membeli Nasi Gerilya melalui GrabFood dalam tiga bulan terakhir? & Pilihan frekuensi & Indikator pembelian ulang dan loyalitas. \\ \hline
        4 & Menu apa yang paling sering dibeli? & Jawaban singkat & Menguji kecocokan dengan menu terlaris TO-GF-04. \\ \hline
        5 & Apa alasan utama memilih Nasi Gerilya? & Pilihan ganda dan jawaban singkat & Membaca daya tarik rasa, porsi, harga, promo, rating, jarak, atau kebiasaan. \\ \hline
        6 & Bagaimana penilaian terhadap rasa, porsi, kesegaran, dan variasi menu? & Skala 1--5 dan catatan & Menguji tema positif TO-GF-07 dan tema kualitas TO-GF-08. \\ \hline
        7 & Bagaimana penilaian terhadap harga, ongkir, voucher, dan promo? & Skala 1--5 dan catatan & Menguji TO-GF-03 dan TO-GF-05 dari sisi pelanggan. \\ \hline
        8 & Apakah harga Nasi Gerilya terasa sesuai dibanding rumah makan Padang lain di GrabFood? & Lebih murah / setara / lebih mahal / tidak tahu & Menguji perbandingan harga dengan kompetitor TO-KP-06. \\ \hline
        9 & Bagaimana penilaian terhadap foto menu, kemasan, dan kondisi makanan saat diterima? & Skala 1--5 dan catatan & Menguji \textit{physical evidence} dan pengalaman produk. \\ \hline
        10 & Apakah pernah mengalami pesanan tidak sesuai, item kurang lengkap, menu habis, atau rasa tidak konsisten? & Ya / Tidak, lalu uraian & Menguji tema keluhan TO-GF-08. \\ \hline
        11 & Jika pernah mengalami masalah, apakah pelanggan menulis ulasan di GrabFood? & Ya / Tidak / Tidak ingat, lalu alasan & Mengurangi risiko \textit{survivorship bias} pada pembacaan ulasan. \\ \hline
        12 & Rumah makan Padang apa yang biasa dibandingkan dengan Nasi Gerilya? & Jawaban singkat & Menguji relevansi Paripurna, Dendeng Batokok, Istana Krakatau, Pondok Gurih, Garuda, atau kompetitor lain. \\ \hline
        13 & Apakah bersedia membeli ulang Nasi Gerilya melalui GrabFood? & Ya / Tidak / Mungkin, lalu alasan & Indikator retensi dan peluang strategi. \\ \hline
        14 & Saran utama untuk meningkatkan pengalaman membeli Nasi Gerilya di GrabFood. & Jawaban terbuka & Masukan strategi produk, harga, promo, proses, atau tampilan toko. \\ \hline
        \bottomrule
    \end{tabularx}
    \tablesource{Diolah peneliti sebagai format survei pelanggan pendukung (2026).}
\end{table}

\begin{table}[!htbp]
    \centering
    \caption{Format Ringkasan Hasil Survei Pelanggan}
    \label{tab:format-ringkasan-survei-pelanggan}
    \footnotesize
    \setlength{\tabcolsep}{4pt}
    \renewcommand{\arraystretch}{0.95}
    \begin{tabularx}{\textwidth}{|Z{3cm}|Y|Y|}
        \toprule
        \textbf{Aspek} & \textbf{Ringkasan Hasil} & \textbf{Implikasi Awal} \\ \hline
        \midrule
        Frekuensi pembelian & Diisi setelah survei & Menguji kekuatan retensi dan pelanggan rutin yang muncul pada ulasan positif. \\ \hline
        Alasan membeli & Diisi setelah survei & Menentukan faktor utama: rasa, porsi, harga, promo, rating, jarak, atau kebiasaan. \\ \hline
        Penilaian produk & Diisi setelah survei & Menguji tema rasa, porsi, kesegaran, dan variasi menu. \\ \hline
        Penilaian harga dan promo & Diisi setelah survei & Menguji persepsi nilai terhadap rentang harga dan voucher. \\ \hline
        Keluhan utama & Diisi setelah survei & Mengonfirmasi keluhan salah item, kelengkapan, stok, dan konsistensi kualitas. \\ \hline
        Perilaku memberi ulasan & Diisi setelah survei & Membantu membaca bias ulasan yang tampil di GrabFood. \\ \hline
        Kompetitor pembanding & Diisi setelah survei & Memvalidasi apakah lima kompetitor awal relevan dari sisi pelanggan atau perlu ditambah dengan kompetitor lain yang lebih sering dibandingkan. \\ \hline
        Niat membeli ulang & Diisi setelah survei & Menjadi dasar strategi retensi dan peningkatan penjualan. \\ \hline
        Saran pelanggan & Diisi setelah survei & Menjadi masukan prioritas strategi. \\ \hline
        \bottomrule
    \end{tabularx}
    \tablesource{Diolah peneliti sebagai format ringkasan survei pelanggan (2026).}
\end{table}
